\begin{statementbox}
\textbf{\textcolor{CStatement}{这条结论解决什么问题？}}

题目没有告诉我们 $f$ 的具体计算式子，只给出一个函数的定义域，
却让我们求 $f(2x+1)$、$f(x+1)$ 或 $f(x^2)$ 的定义域。
\textbf{不用猜出 $f$ 的式子，也可以判断哪些 $x$ 能代入。}

\medskip

\begin{center}\begin{tikzpicture}[font=\small,>=stealth,line join=round]
\draw[gray!45,rounded corners=3pt] (0,0) rectangle (14.2,9.4);
\node at (7.1,8.8) {把允许的输入范围拉回到自变量轴};

\begin{scope}[shift={(2.2,2.6)},scale=1.35]
\fill[blue!8] (-.3,1) rectangle (2.6,3);
\draw[->] (-.5,0)--(2.8,0) node[right]{$x$};\draw[->] (0,-.7)--(0,3.6) node[above]{$t$};
\draw[teal!80!black,thick,domain=.2:2.2] plot (\x,{2*\x-1}) node[above right]{$t=2x-1$};
\foreach \x/\y in {1/1,2/3}{\draw[dashed,gray] (\x,0)--(\x,\y)--(0,\y);\fill (\x,\y) circle(1.5pt);\node[below] at (\x,0){$\x$};\node[left] at (0,\y){$\y$};}
\draw[blue,line width=2pt] (0,1)--(0,3);\draw[orange!90!black,line width=2pt] (1,0)--(2,0);
\end{scope}
\node[align=left,anchor=west] at (8,5.1) {外层允许：$t\in[1,3]$\\[8pt] 内层关系：$t=2x-1$\\[8pt] $1\le2x-1\le3$\\[8pt] 自变量：$x\in[1,2]$};
\node[text width=13.2cm,align=center] at (7.1,.65) {先检查括号里的数，再解出 $x$ 的范围。};
\end{tikzpicture}
\end{center}
\textbf{\textcolor{CStatement}{先认清两个数}}
\begin{itemize}[leftmargin=*,itemsep=0.3em]
  \item $x$：题目要求我们找范围的数，也是 $f[g(x)]$ 的自变量。
  \item $g(x)$：先用 $x$ 算出来，再放进 $f$ 括号里的数。
\end{itemize}

\end{statementbox}
\clearpage
\begin{statementbox}
\textbf{\textcolor{CStatement}{结论一：知道 $f$ 能接收哪些数，求 $f[g(x)]$ 的定义域}}

设 $f(t)$ 的定义域为 $D$。要使 $f[g(x)]$ 有意义，必须先能算出 $g(x)$，
并且算出的数属于 $D$。所以
\[
 \boxed{\ f[g(x)]\text{的定义域}=\{x\mid g(x)\text{有意义且 }g(x)\in D\}\ }.
\]
如果 $g(x)$ 对每个实数 $x$ 都有意义，就直接解 $g(x)\in D$。

\medskip
\textbf{\textcolor{CStatement}{结论二：知道 $f(ax+b)$ 的定义域，反求 $f(x)$ 的定义域}}

如果 $a\ne0$，而且题目给出的是 $f(ax+b)$ 的\textbf{完整定义域} $A$，
那么先让 $x$ 在 $A$ 中取值，算出 $ax+b$ 能得到的所有数，这些数恰好就是 $f$ 能接收的数：
\[
 \boxed{\ f(t)\text{的定义域}=\{ax+b\mid x\in A\},\quad a\ne0\ }.
\]
例如，$f(2x-1)$ 的定义域是 $[1,4]$，则 $f(t)$ 的定义域是 $[1,7]$。

\medskip
\textbf{\textcolor{CStatement}{一定要看清的限制}}

第二条\textbf{不能不加判断地换成任意 $g(x)$}。
例如 $x^2$ 永远不可能是负数：知道 $f(x^2)$ 的定义域，
仍可能不知道 $f$ 能否接收负数。
因此，只有当 $x$ 在括号内式子本身允许的全部范围中变化时，该式子\textbf{能够取遍所有实数}，
才能仅凭复合函数的完整定义域，确定 $f$ 的完整定义域。
\end{statementbox}
